\documentclass[a4paper,twoside]{article}

\usepackage{morewrites}

\usepackage[svgnames,dvipsnames,x11names]{xcolor}
\usepackage{graphicx}
\usepackage[english]{babel}
\usepackage[colorlinks=true, linkcolor=NavyBlue, urlcolor=NavyBlue, citecolor=NavyBlue]{hyperref}
\usepackage[a4paper]{geometry}
\usepackage{ts-fonts}
\usepackage{float}
\usepackage{seqsplit}

\usepackage{lastpage}
\usepackage{refcount}
\makeatletter
\def\ps@plain{
  \let\@oddhead\@empty
  \let\@evenhead\@empty
  \def\@oddfoot{
    \hfil
    \ifnum\getpagerefnumber{LastPage}>1\relax
      \thepage
    \fi
    \hfil}
  \let\@evenfoot\@oddfoot
}
\makeatother

\title{Partials \& Override Precedence}
\author{}\date{}

\begin{document}

\maketitle

\thispagestyle{plain}
TeXSmith renders \LaTeX{} through reusable partials (Jinja templates) for list items, code blocks, images, etc. The override order is \textbf{explicit and enforced}:

\begin{enumerate}
\item \textbf{Template overrides} --- declared under \texttt{latex.template.override} in \texttt{manifest.toml}. These always win.
\item \textbf{Fragment overrides} --- declared via \texttt{partials} in \texttt{fragment.toml} or a fragment entrypoint. They apply after core defaults but are superseded by template overrides.
\item \textbf{Core defaults} --- shipped with TeXSmith (\texttt{src/texsmith/adapters/latex/partials}).

\end{enumerate}
Conflicts \& requirements:
- Two fragments cannot override the same partial; TeXSmith raises a \texttt{TemplateError}.
- Templates and fragments can declare \texttt{required\_partials} (list of names without extensions). Missing entries abort the render with a clear error.
- Providers are tracked so diagnostics can point to the template or fragment responsible.

Usage:
- Templates: list override paths in \texttt{manifest.toml} under \texttt{latex.template.override} (relative to the template root, \texttt{template/overrides}, or \texttt{overrides/}).
- Fragments: add \texttt{partials = ["strong.tex", "codeblock.tex"]} (or a \texttt{\{ name = "path" \}} mapping) plus optional \texttt{required\_partials} inside \texttt{fragment.toml}. Entry points may return the same fields on the \texttt{Fragment} object.

Guidance:
- Prefer template overrides for document-wide styling or publisher branding.
- Use fragment overrides to co-locate rendering tweaks with the feature they own (callouts, code, glossary).
- Keep a single source of truth per partial; if a fragment needs to opt out, expose a fragment attribute instead of overlapping overrides.

\end{document}
